\chapter{I LUOGHI FREDDI DELL'UNIVERSO}\label{i-luoghi-freddi-dell-universo}

\hfill\break

Un venerdì sera, tornai a casa dopo una riunione alla Perielio, aprii
con le chiavi la porta di casa e trovai Molly seduta alla tastiera del
mio PC.

La scrivania era nell'angolo a sud-ovest del soggiorno a ridosso di una
finestra e dava le spalle alla porta. Molly si voltò di mezzo giro e mi
rivolse uno sguardo sorpreso. Al tempo stesso, abilmente, cliccò su
un'icona e uscì dal programma che stava usando.

«Molly?»

Non ero sorpreso di trovarla lì. Moll passava quasi tutti i fine
settimana con me; aveva un duplicato delle chiavi. Non aveva però mai
mostrato alcun interesse per il mio PC.

«Non hai chiamato» disse lei.

Ero stato a un incontro con i rappresentanti dell'assicurazione che
sottoscriveva la copertura per i dipendenti della Perielio. Mi avevano
detto di prevedere una riunione di due ore, ma alla fine si era rivelato
un aggiornamento di venti minuti sulla fatturazione e, quando era
terminato, avevo pensato che avrei fatto prima a prendere l'auto per
tornare a casa e che forse sarei persino arrivato prima di Molly, se si
fosse fermata a prendere il vino.

Il modo in cui mi fissò a lungo ebbe un effetto tale da farmi sentire
obbligato a spiegarle tutto prima di chiederle cosa stesse facendo con i
miei file.

Mentre percorrevo la stanza si mise a ridere, una di quelle risate
imbarazzate, dispiaciute: ``Ma guarda cosa mi hai beccato a fare.'' La
mano destra indugiava sul touchpad del PC. Si voltò di nuovo verso il
monitor. Sullo schermo, il cursore si precipitò verso l'icona di
arresto.

«Aspetta» dissi.

«Che c'è, devi usarlo?»

Il cursore puntava sul bersaglio. Misi la mano su quella di Molly. «In
realtà, vorrei sapere cosa stavi facendo.»

Era tesa. Una vena le pulsava sulla pelle rosa proprio davanti
all'orecchio. «Ho fatto come se fossi a casa mia. Uhm, forse un po'
troppo? Non credevo ti sarebbe dispiaciuto.»

«Dispiaciuto di cosa, Moll?»

«Dispiaciuto se avessi usato il tuo computer.»

«Usato per cosa?»

«Niente di importante. Davo solo un'occhiata.»

Ma non poteva essere il computer a incuriosirla. Aveva cinque anni, era
quasi un pezzo d'antiquariato. Al lavoro usava dispositivi più
sofisticati. E avevo riconosciuto il programma che si era precipitata a
chiudere quando ero entrato. Era un programma di gestione familiare,
quello che usavo per pagare le bollette, far quadrare i conti e gestire
i miei contatti.

«Mi sembrava un foglio di calcolo» dissi.

«Ci sono finita per caso. Il tuo desktop mi ha confuso. Sai, ognuno
organizza le cose in modo diverso. Mi spiace, Tyler. Credo di essere
stata invadente.» Sfilò di scatto la mano da sotto la mia e cliccò sul
comando di arresto. Il desktop si rimpicciolì e sentii la ventola del
processore sibilare finché non si spense. Molly si alzò, lisciandosi la
camicetta. Si dava sempre una sistemata ai vestiti quando si alzava in
piedi. Per mettere le cose in ordine. «Che ne dici se comincio a
preparare la cena?» Si voltò dandomi le spalle e si diresse verso la
cucina.

La vidi svanire dietro le porte a battente. Contai fino a dieci e le
andai dietro.

Stava prendendo le padelle dallo scaffale a muro. Mi lanciò un'occhiata
e distolse lo sguardo.

«Molly, se vuoi sapere qualcosa non devi fare altro che chiedermelo.»

«Ah, tutto qui? Va bene.»

«Molly\ldots»

Poggiò la padella sul fornello con una cura esagerata, quasi fosse
fragile. «Devo scusarmi ancora? Okay, Tyler, mi spiace di aver usato il
tuo computer senza permesso.»

«Non ti sto accusando di niente, Moll.»

«Allora perché ne stiamo parlando? Perché ho come l'impressione che
passeremo tutto il resto della serata a parlare di questo?» Le si
inumidirono gli occhi. Le lenti colorate divennero di una tonalità più
scura di verde smeraldo. «Ero solo un po' curiosa.»

«Curiosa di che, delle mie bollette?»

«Di \emph{te}.» Trascinò una sedia dal tavolo della cucina. La gamba
della sedia si incastrò in quella del tavolo e Molly per liberarla le
diede uno strattone. Si sedette e incrociò le braccia. «Sì, forse anche
delle cose insignificanti. Forse \emph{soprattutto} delle cose
insignificanti.» Chiuse gli occhi e scosse la testa. «Mentre lo dico mi
sembra di essere una specie di stalker. Ma sì, delle bollette, della
marca di dentifricio, del tuo numero di scarpe, sì. Sì, vorrei essere
qualcosa di più di una scopata del fine settimana. Lo ammetto.»

«Non c'è bisogno di frugare nei miei file per questo.»

«Forse non avrei dovuto, se\ldots»

«Se?»

Scosse la testa. «Non voglio litigare.»

«A volte è meglio finire quello che si inizia.»

«Be', se tu non facessi così, per esempio. Ogni volta che ti senti
minacciato, fai il \emph{distaccato}. Diventi tutto freddo e riservato e
analitico, come se io fossi un documentario sulla natura che stai
guardando in TV. Lo schermo sembra abbassarsi, ma rimane \emph{sempre}
lì, vero? Tutto il mondo resta dall'altra parte. Ecco perché non parli
mai di te. Ecco perché ho passato un anno ad aspettare che notassi che
ero più di un mobile da ufficio. Quello sguardo fisso stupido infinito
freddo, che guarda la vita come se fosse il telegiornale serale, come se
fosse qualche misera guerra dall'altra parte del pianeta dove tutte le
persone hanno nomi impronunciabili.»

«Molly\ldots»

«Lo so che siamo tutti scoppiati, Tyler, tutti noi che siamo nati nello
Spin. Disturbo da stress pretraumatico, o com'era che ci hai chiamati?
Una generazione di tipi grotteschi. Ecco perché siamo tutti divorziati o
promiscui o iperreligiosi o depressi o maniaci o \emph{freddi}. Abbiamo
tutti un'ottima scusa per i nostri comportamenti sbagliati, compresa me,
e se essere questo grande esempio di disponibilità studiata nei dettagli
ti fa andare avanti, va bene, lo capisco. Ma volere di più dev'essere
giusto anche per me. È giusto, infatti, è umano che io desideri
toccarti. Non solo scoparti. Toccarti.»

Disse tutte queste cose e poi, rendendosi conto di aver finito, spiegò
le braccia e aspettò una mia reazione.

Pensai di farle anch'io un discorso. Le avrei detto che \emph{sì},
provavo una forte passione per lei. Magari non l'avevo dato a vedere, ma
mi ero accorto di lei fin dai tempi in cui ero andato a lavorare alla
Perielio. Accorto delle linee e della dinamica del suo corpo, di come si
alzava, camminava, si stiracchiava o sbadigliava; accorto del suo
guardaroba pastello e della farfalla di bigiotteria che indossava su una
catenina d'argento; accorto dei suoi stati d'animo, degli impulsi e del
suo inventario di sorrisi, bronci e gesti. Quando chiudevo gli occhi
vedevo il suo viso e quando andavo a dormire era ciò che guardavo. Mi
piacevano il suo aspetto e la sua essenza: il sapore salato della sua
lingua, la cadenza della sua voce, l'arco delle sue dita e le parole che
scrivevano sul mio corpo.

Pensai a tutto questo, ma non mi riuscì di dirglielo.

Non era proprio una bugia. Ma non era nemmeno la verità.

Alla fine facemmo la pace scambiandoci qualche carineria, brevi pianti e
abbracci conciliatori, lasciammo perdere la discussione e io mi prestai
a fare l'aiuto cuoco mentre lei preparava un sugo per la pasta veramente
buono. La tensione cominciò ad allentarsi, ci coccolammo per un'ora
guardando notiziari fino a mezzanotte (aumento della disoccupazione, un
dibattito elettorale, qualche misera guerra dall'altra parte del
pianeta) e poi fummo pronti per andare a letto. Molly spense la luce
prima di fare l'amore, la camera era buia e la finestra aperta e il
cielo uniforme e vuoto. Quando venne, inarcò la schiena e gemette. Il
respiro era dolce e profumato di latte. Divisi, ma continuando a
toccarci con la mano sulla coscia, ci scambiammo frasi incompiute.
Dissi: «Sai, la \emph{passione}» e lei: «A letto, sicuramente.»

Si addormentò subito. Io un'ora dopo ero ancora sveglio.

Scesi piano dal letto e non registrai alcun cambiamento nel ritmo del
suo respiro. Mi infilai un paio di jeans e lasciai la camera. Nelle
notti insonni come quella, un po' di Drambuie di solito mi aiutava a far
tacere l'irritante monologo interiore, le istanze presentate dal dubbio
all'affaticato prosencefalo. Ma prima di andare in cucina, mi sedetti al
PC e richiamai il mio programma di gestione familiare.

Non c'era nulla che rivelasse ciò che Moll stava guardando. Difatti non
c'era niente fuori posto, stando a quello che vedevo. Tutti i nomi e i
numeri sembravano uguali a prima. Forse aveva trovato qualcosa che
l'aveva fatta sentire più vicina a me. Se era davvero ciò che voleva.

O forse era stata una ricerca inutile. Forse non aveva trovato proprio
niente.

\separatorefiori

Nelle settimane precedenti alle elezioni di novembre vidi Jason più
spesso. La sua malattia stava diventando più attiva nonostante dosi
maggiori di farmaco, probabilmente a causa dello stress provocato dal
contrasto crescente con suo padre. (E.D. aveva annunciato l'intenzione
di ``riprendersi'' la Perielio per liberarla da quella che considerava
una cricca di burocrati arrampicatori e scienziati allineati con Wun Ngo
Wen - una minaccia vuota, secondo Jason, ma potenzialmente distruttiva e
imbarazzante.)

Jase mi teneva accanto a sé nel caso fosse necessario somministrargli
degli antispasmodici in qualche momento critico, cosa che ero disposto a
fare entro i limiti della legge e dell'etica professionale. Mantenere
Jase operativo a breve termine era il massimo che la scienza medica
poteva fare per lui, e rimanere operativo il tempo sufficiente per
battere con abilità E.D. Lawton era, in quel momento, tutto ciò che
importava a Jase.

E così passavo molto tempo nell'ala VIP della Perielio, di solito con
Jason ma spesso con Wun Ngo Wen. Questo mi rendeva oggetto di diffidenza
da parte degli addetti di Wun, un assortimento di autorità
subgovernative (giovani rappresentanti del Dipartimento di Stato, della
Casa Bianca, della Difesa nazionale, del Comando spaziale, eccetera) e
accademici assoldati per tradurre, studiare e classificare i cosiddetti
archivi marziani. Agli occhi di queste persone la mia libertà di
avvicinare Wun era contro le regole e sgradita. Ero un mercenario. Una
nullità. Ma era per quel motivo che Wun preferiva la mia compagnia: non
avevo secondi fini da promuovere o proteggere. E siccome insisteva, di
tanto in tanto dei burberi lacchè mi accompagnavano attraverso le varie
porte che separavano l'alloggio climatizzato dell'ambasciatore marziano
dal caldo della Florida e dal vasto mondo esterno.

In una di quelle occasioni trovai Wun Ngo Wen seduto sulla sua sedia di
vimini - qualcuno aveva procurato un poggiapiedi coordinato in modo che
i piedi non gli penzolassero - che fissava pensieroso il contenuto di
una provetta di vetro. Gli chiesi cosa contenesse.

«Replicatori» rispose.

Indossava un abito che sembrava essere stato confezionato per un
dodicenne tarchiato: aveva fatto una presentazione per una delegazione
del Congresso. Malgrado l'esistenza di Wun non fosse stata annunciata
formalmente, nelle ultime settimane c'era stato un viavai continuo di
visitatori, sia stranieri che locali, accettati dalla sicurezza.
L'annuncio ufficiale sarebbe stato fatto dalla Casa Bianca subito dopo
le elezioni e in seguito Wun sarebbe stato davvero molto impegnato.

Guardai la provetta da una buona distanza di sicurezza.
\emph{Replicatori}. Mangiatori di ghiaccio. Semi di una biologia
inorganica.

Wun sorrise. «Ti fa paura? Non averne, ti prego. Ti assicuro che il
contenuto è completamente inattivo. Pensavo che Jason te l'avesse
spiegato.»

Lo aveva fatto. Un po'. Dissi: «Sono dispositivi microscopici.
Semiorganici. Si riproducono in condizioni di freddo estremo e vuoto.»

«Sì, bene, nella sostanza è corretto. E Jason ti ha spiegato il loro
scopo?»

«Andare a popolare la galassia. Per inviarci dati.»

Wun annuì lentamente, come se anche questa risposta fosse
sostanzialmente corretta ma meno soddisfacente. «Questo è il manufatto
tecnologico più sofisticato che abbiano prodotto le Cinque Repubbliche,
Tyler. Non avremmo mai potuto sostenere il tipo di attività industriale
che la tua gente esercita su scala tanto allarmante - transatlantici,
uomini sulla Luna, città enormi\ldots»

«Da ciò che ho visto, le vostre città sono alquanto imponenti.»

«Solo perché le costruiamo in un gradiente gravitazionale più lieve.
Sulla Terra quelle torri si sarebbero sbriciolate sotto il loro stesso
peso. Ma io penso che questo, il contenuto di questa provetta, sia il
nostro equivalente di trionfo ingegneristico, qualcosa di così complesso
e difficile da realizzare che ne siamo, forse a ragione, orgogliosi.»

«Sono certo che lo siate.»

«Allora vieni qui, così potrai apprezzarlo. Non aver paura.» Mi fece
cenno di avvicinarmi e attraversai la stanza per sedermi su una sedia di
fronte a lui. Immagino che, da lontano, avremmo potuto essere scambiati
per due amici qualsiasi che discutono di qualcosa. I miei occhi, però,
non si staccavano dalla provetta. Me la porse, me la offrì. «Su,
prendila» mi esortò.

Presi la fialetta tra pollice e indice e la sollevai per fare in modo
che la luce del soffitto la attraversasse. Il contenuto sembrava acqua
comune con una lucentezza lievemente oleosa. Tutto lì.

«Per apprezzarla veramente,» disse Wun, «devi capire cos'hai in mano. In
quella provetta, Tyler, ci sono circa trenta o quarantamila singole
cellule artificiali in una sospensione di glicerina. Ogni cellula è una
ghianda.»

«Ti intendi di ghiande?»

«Ho cominciato a leggere qualcosa sull'argomento. È una metafora comune.
Ghiande e querce, giusto? Quando tieni in mano una ghianda, stai
reggendo un potenziale albero di quercia, e non solo una singola quercia
ma tutta la progenie di quella quercia per secoli e secoli. Legno di
quercia sufficiente a costruire intere città\ldots{} le città sono fatte
di quercia?»

«No, ma non importa.»

«Quella che hai in mano è una ghianda. Del tutto dormiente, come ti
dicevo, e in effetti \emph{quel} campione in particolare potrebbe essere
morto, considerando il tempo che ha passato alla temperatura ambiente
terrestre. Se la analizzi potresti trovare al massimo qualche traccia
insolita di sostanza chimica.»

«Ma?»

«Ma\ldots{} se la metti in un ambiente freddo, privo d'aria, ghiacciato,
un ambiente come la Nube di Oort, vedrai che prende vita, Tyler! Inizia,
molto lentamente ma con molta pazienza, a crescere e riprodursi.»

La Nube di Oort. Sapevo della Nube di Oort dalle conversazioni con Jason
e dai romanzi speculativi che ancora leggevo di tanto in tanto. Era una
schiera nebulosa di corpi cometari le cui dimensioni andavano
all'incirca dall'orbita di Plutone alla metà della sua distanza con la
stella più vicina. Quei piccoli corpi non erano per niente ammassati -
occupavano un enorme volume di spazio, quasi inconcepibile - ma la loro
massa totale equivaleva a venti, trenta volte quella della Terra,
soprattutto sotto forma di ghiaccio sporco.

Un sacco di roba da mangiare, se mangi ghiaccio e polvere.

Wun si piegò in avanti sulla sedia. I suoi occhi, acquattati nella pelle
di cuoio raggrinzito, brillavano. Sorrise, gesto che avevo imparato a
interpretare come segno di sincerità: i marziani sorridono quando
parlano col cuore.

«Questo, tra la mia gente, ha provocato non poche controversie. Ciò che
reggi in mano ha il potere di trasformare profondamente non solo il
nostro sistema solare, ma molti altri. E ovviamente il risultato è
incerto. I replicatori sono inorganici nel senso comune del termine, ma
allo stesso tempo sono \emph{vivi}. Vivono cicli di retroazione
autocatalitici, soggetti a modifiche da parte della pressione
ambientale. Proprio come gli esseri umani o i batteri o, o\ldots»

«O i \emph{murkud}» dissi.

Fece un gran sorriso. «O i \emph{murkud}.»

«In altre parole, potrebbero evolversi.»

«Si \emph{evolveranno} certamente e in maniera imprevedibile. Ma abbiamo
posto dei limiti. O almeno crediamo di averlo fatto. Come ti dicevo, le
polemiche abbondano.»

Ogni volta che Wun mi parlava di politica marziana, mi immaginavo uomini
e donne grinzosi su podi in acciaio inossidabile, avvolti in toghe color
pastello, che discutevano di argomenti astratti. In effetti, invece,
insisteva Wun, i parlamentari marziani si comportavano più come
agricoltori squattrinati che bisticciano a un'asta di grano; e i
vestiti\ldots{} be', non provavo nemmeno a immaginarmeli i vestiti;
nelle occasioni formali i marziani di entrambi i sessi tendevano a
vestirsi come la regina di cuori in un mazzo di carte Bicycle.

Ma se le discussioni erano state lunghe e accorate, il progetto in sé fu
invece relativamente semplice. I replicatori sarebbero stati lanciati a
casaccio verso i confini lontani e freddi del sistema solare. Qualche
frazione infinitesimale di quei replicatori si sarebbe posata su due o
tre nuclei cometari che costituiscono la Nube di Oort. Lì avrebbero
cominciato a riprodursi.

Le loro informazioni genetiche, disse Wun, erano codificate in molecole
termicamente instabili ovunque più calde delle lune di Nettuno. Ma
nell'ambiente iperfreddo per cui erano stati progettati, i filamenti
submicroscopici dei replicatori avrebbero iniziato un metabolismo lento
e accurato. Erano progettati per crescere a una lentezza tale che al
confronto lo sviluppo di un pino dal cono setoloso sarebbe parso
fulmineo, eppure sarebbero cresciuti, assorbendo tracce di sostanze
volatili e molecole organiche e modellando il ghiaccio in pareti
cellulari, costole, nervature e giunture.

Nel momento in cui i replicatori avessero consumato all'incirca qualche
centinaio di metri cubi di nucleo cometario, le loro interconnessioni
avrebbero iniziato a diventare più complesse e il loro comportamento più
risoluto. Si sarebbero formate appendici altamente sofisticate, occhi di
ghiaccio e carbonio per scrutare l'oscurità stellare.

In circa un decennio la colonia di replicatori sarebbe diventata
un'entità di gruppo sofisticata in grado di registrare e trasmettere
dati rudimentali sul suo ambiente. Avrebbe guardato il cielo e chiesto:
«C'è un corpo oscuro di dimensioni planetarie che gira intorno alla
stella più vicina?»

Porre la domanda e ottenere una risposta avrebbe richiesto altri
decenni, ma almeno inizialmente il risultato sarebbe stato scontato: sì,
i due mondi che giravano intorno a quella stella erano corpi oscuri,
Terra e Marte.

Tuttavia - lentamente, con pazienza e ostinazione - i replicatori
avrebbero confrontato quei dati e li avrebbero ritrasmessi al loro punto
di origine: a noi o comunque ai nostri satelliti in ascolto.

Poi, con la sua senescenza da macchina complessa, la colonia di
replicatori si sarebbe suddivisa in singoli ammassi di cellule semplici,
avrebbe identificato un'altra stella luminosa o vicina e utilizzato le
sostanze volatili accumulate estratte dal nucleo cometario ospite per
spingere i semi oltre il sistema solare. (Avrebbero lasciato un loro
minuscolo frammento che avrebbe agito da ripetitore radio, un nodo
passivo in una rete in crescita.)

Questi semi di seconda generazione avrebbero vagato nello spazio
interstellare per anni, decenni, millenni. Alla fine la maggior parte
sarebbe andata distrutta, persa in traiettorie infruttuose o trascinata
in vortici gravitazionali. Alcuni, incapaci di sfuggire all'attrazione
debole e lontana del Sole, sarebbero ricaduti nella Nube di Oort e
avrebbero ripetuto il processo, mangiando ghiaccio in modo stupido ma
paziente e registrando informazioni ridondanti. Se si fossero incontrati
due ceppi, si sarebbero scambiati materiale cellulare, avrebbero ridotto
gli errori di duplicazione indotti dal tempo o dalle radiazioni e
prodotto una discendenza quasi simile ma non esattamente uguale a loro.

Alcuni avrebbero raggiunto l'alone ghiacciato di una stella vicina e
ricominciato da capo il ciclo, questa volta raccogliendo informazioni
nuove, che infine ci avrebbero inviato come esplosioni di dati, brevi
orgasmi digitali. «Stella binaria,» forse avrebbero detto, «nessun corpo
planetario oscuro»; oppure avrebbero detto: «Nana bianca, un corpo
planetario oscuro.»

E il ciclo si sarebbe ripetuto ancora.

E ancora.

E ancora, da una stella all'altra, in sequenza, secoli dopo millenni,
atrocemente lento, ma abbastanza veloce per i tempi della galassia -
mentre noi cronometravamo l'universo esterno dalla nostra bara. I nostri
giorni sarebbero stati centinaia di migliaia dei loro anni e un decennio
del nostro tempo lento li avrebbe visti invadere quasi tutta la
galassia.

Le informazioni passate da nodo a nodo alla velocità della luce
sarebbero state inoltrate, avrebbero modificato il comportamento,
avrebbero indirizzato nuovi replicatori verso territori inesplorati,
avrebbero soppresso le informazioni ridondanti in modo che i nodi
centrali non fossero sovraccaricati. In pratica, avremmo collegato la
galassia per una specie di pensiero rudimentale. I replicatori avrebbero
costruito una rete neurale grande quanto il cielo notturno che avrebbe
comunicato con noi.

C'erano dei rischi? Certamente ce n'erano.

In assenza dello Spin, disse Wun, i marziani non avrebbero mai approvato
una simile appropriazione arrogante delle risorse della galassia. Non
era un semplice atto di esplorazione; era un \emph{intervento}, un
riordino imperiale dell'ecologia galattica. Se là fuori ci fossero state
altre specie senzienti - e l'esistenza degli Ipotetici aveva pressoché
risposto in maniera affermativa a quella domanda - lo spargimento dei
replicatori forse sarebbe stato interpretato come un'aggressione che
avrebbe potuto provocare ritorsioni.

I marziani avevano riconsiderato questo rischio soltanto quando avevano
scoperto sui loro poli nord e sud le strutture in costruzione dello
Spin.

«Lo Spin rende irrilevanti le obiezioni,» disse Wun, «o quasi. Con un
po' di fortuna i replicatori ci diranno qualcosa di importante sugli
Ipotetici o almeno sull'entità delle loro opere nella galassia. Potremmo
riuscire a comprendere lo scopo dello Spin. In caso contrario, i
replicatori serviranno come una sorta di segnale d'avvertimento per
altre specie intelligenti che si troveranno ad affrontare lo stesso
problema. Un'analisi accurata suggerirebbe a un osservatore premuroso lo
scopo per cui è stata costruita la rete. Altre civiltà potrebbero
scegliere di sfruttarla. La conoscenza potrebbe aiutarli a proteggersi.
Per riuscire dove noi abbiamo fallito.»

«Pensi che falliremo?»

Wun scrollò le spalle. «Non abbiamo già fallito? Il Sole ormai è
vecchissimo. Lo sai, Tyler. Niente dura all'infinito. E date le
circostanze, per noi anche ``all'infinito'' non è un tempo lunghissimo.»

Forse fu il modo in cui lo disse, con il suo triste e sincero sorrisetto
da marziano, piegandosi in avanti sulla sedia di vimini, ma il peso di
quella affermazione fu sommessamente scioccante.

Non ne fui sorpreso. Tutti sapevamo di essere condannati. Condannati, a
dir poco, a passare l'intera vita sotto un guscio che era l'unica cosa a
proteggerci da un sistema solare ostile. La luce solare che aveva reso
Marte abitabile avrebbe cotto la Terra se la membrana dello Spin fosse
stata rimossa. E persino Marte (nel suo stesso involucro oscuro) stava
scivolando via rapidamente dalla cosiddetta zona abitabile. La stella
mortale, madre di ogni vita, era entrata nella senescenza sanguinaria e
ci avrebbe uccisi senza rimorsi.

La vita era nata ai margini di una reazione nucleare instabile. Era vero
ed era sempre stato vero; era stato vero prima dello Spin, anche quando
il cielo era sgombro e le notti estive brillavano di stelle distanti, di
nessun interesse. Era stato vero ma non aveva avuto importanza perché la
vita umana era breve; innumerevoli generazioni sarebbero vissute e morte
nel tempo di un battito solare. Ma in quel momento, grazie a Dio,
stavamo sopravvivendo al Sole. O saremmo finiti per essere cenere e
ruotare intorno al suo cadavere o saremmo stati preservati nella notte
eterna, ninnoli incapsulati senza una vera casa nell'universo.

«Tyler? Va tutto bene?»

«Sì» dissi. Per qualche ragione, stavo pensando a Diane. «Forse la cosa
migliore che possiamo sperare è capirci qualcosa prima che cali il
sipario.»

«Sipario?»

«Prima della fine.»

«Non è una gran consolazione» ammise Wun. «Comunque sì, può essere la
cosa migliore in cui sperare.»

«La tua gente sapeva dello Spin da millenni. E in tutto quel tempo non
siete riusciti a sapere niente sugli Ipotetici?»

«No. Mi dispiace. Questo non sono in grado di dirvelo. Sulla natura
fisica dello Spin abbiamo solo qualche ipotesi.» (Che Jason di recente
aveva tentato di spiegarmi: qualcosa sui quanti temporali, per lo più
matematica e ben oltre la portata dell'ingegneria applicata, marziana o
terrestre.) «Sugli Ipotetici, niente di niente. Riguardo a quello che
vogliono da noi\ldots» Scrollò le spalle. «Solo altre ipotesi. La
domanda che ci siamo posti è stata: cosa aveva di speciale la Terra per
essere incapsulata? Perché gli Ipotetici hanno aspettato a coprire Marte
con uno Spin e cos'è che li ha spinti a scegliere quel particolare
momento della nostra storia?»

«Hai delle risposte?»

Uno dei suoi addetti bussò alla porta e l'aprì. Un tizio quasi calvo con
un abito nero su misura. Parlava con Wun ma guardava me: «Volevo solo
ricordarle che sta arrivando il rappresentante dell'UE. Fra cinque
minuti.» Tenne la porta aperta, in attesa. Mi alzai.

«Alla prossima» disse Wun.

«Presto, spero.»

«Appena riesco a organizzarmi.»

Era tardi e per quel giorno avevo finito. Uscii dalla porta nord. Sulla
strada per il parcheggio mi fermai vicino alla staccionata di legno dove
stavano costruendo la nuova ala della Perielio. Tra i buchi del muro di
sicurezza riuscii a vedere un comune edificio a blocchi di cemento,
enormi serbatoi a pressione esterna, tubi spessi come barili che
attraversavano strombature in calcestruzzo. Il terreno era disseminato
di materiale isolante giallo in PTFE e tubi di rame avvolti. Un
caposquadra con un casco bianco urlava ordini a uomini con occhiali di
protezione e scarponi con la punta d'acciaio, che spingevano carriole.
Uomini che costruivano un'incubatrice per un nuovo genere di vita.

Era lì, in culle di elio liquido, che i replicatori sarebbero stati
allevati e preparati per il loro lancio nei luoghi freddi dell'universo:
i nostri eredi, in un certo senso, destinati a vivere più a lungo e a
viaggiare più lontano di quanto gli esseri umani avrebbero mai fatto. Il
nostro dialogo finale con l'universo. Sempre che E.D. non la spuntasse
annullando tutto il progetto.

\separatorefiori

Quel fine settimana io e Molly facemmo una passeggiata sulla spiaggia.

Era un sabato sereno di fine ottobre. Avevamo percorso circa cinquecento
metri di sabbia zeppa di mozziconi di sigaretta prima che la giornata si
facesse esageratamente calda e il Sole diventasse persistente. L'oceano
rifletteva la luce restituendo puntini abbaglianti, come se banchi di
diamanti nuotassero al largo. Molly indossava pantaloncini e sandali,
una maglietta bianca di cotone che le aderiva al corpo in maniera
seducente e un berretto con la visiera tirata giù per farsi ombra sugli
occhi.

«Non l'ho mai capito» disse, passandosi il polso sulla fronte e
voltandosi per guardare le proprie orme sulla sabbia.

«Cosa, Moll?»

«Il Sole. Cioè, la luce solare. Questa luce. È finta, lo dicono tutti
ma, mio Dio, che caldo: il caldo è vero.»

«Il Sole non è del tutto finto. Questo \emph{sole} che vediamo non è il
vero Sole, ma questa luce ha avuto origine lì. È gestito dagli
Ipotetici, le lunghezze d'onda ridotte e filtrate\ldots»

«Lo so, ma intendo dire il modo in cui si sposta nel cielo. Alba,
tramonto. Se è solo una proiezione, come mai sembra lo stesso in Canada
e Sud America? Se la barriera dello Spin è solo a qualche centinaio di
chilometri dalla superficie del nostro pianeta?»

Le dissi cosa mi aveva detto Jason una volta: il \emph{sole} finto non
era un'illusione proiettata su uno schermo, era una replica gestita
della luce solare che passava \emph{attraverso} lo schermo da una fonte
lontana all'incirca centocinquanta milioni di chilometri, come un
software di ray tracing in scala colossale.

«Un trucco di scena piuttosto elaborato, cazzo» disse Molly.

«Se avessero agito diversamente, saremmo morti tutti anni fa. L'ecologia
planetaria ha bisogno di un giorno di ventiquattro ore.» Avevamo già
perso un certo numero di specie che dipendeva dalla luce della Luna per
nutrirsi e accoppiarsi.

«Ma è una bugia.»

«Se vuoi definirla così\ldots»

«Una bugia, sì, la chiamo bugia. Sto qui con la luce di una bugia in
faccia. Una bugia da cui può venirti il cancro alla pelle. Eppure
continuo a non capire. Immagino che non ci capiremo niente finché non
capiremo gli Ipotetici. Se mai accadrà. Ne dubito.»

«Una bugia non la si capisce,» continuò Molly mentre costeggiavamo
un'antica passerella diventata bianca per il sale, «fino a quando non si
comprende la motivazione che c'è dietro.» Lo disse guardandomi di
sottecchi, gli occhi nascosti sotto il berretto, mandandomi messaggi che
non riuscivo a decifrare.

Passammo il resto del pomeriggio nel mio appartamento climatizzato a
leggere e ascoltare musica, ma Moll era inquieta e io non avevo
accettato del tutto la sua incursione nel mio computer, altro evento
indecifrabile. Amavo Molly. O almeno mi dicevo di amarla. O comunque, se
quello che provavo per lei non era amore, ne era quanto meno
un'imitazione plausibile, un surrogato convincente.

Di lei mi preoccupava il fatto che fosse sempre molto imprevedibile,
ossessionata dallo Spin come tutti noi. Non riuscivo a comprarle dei
regali: c'erano cose che voleva ma, a meno che non avesse espresso a
voce il proprio apprezzamento davanti alla vetrina di un negozio, non
riuscivo a immaginare cosa fossero. Teneva ben nascoste le sue esigenze
più profonde. Forse, come la maggior parte delle persone riservate, dava
per scontato che anch'io tenessi per me dei segreti importanti.

Avevamo appena finito di cenare e stavamo cominciando a mettere in
ordine quando squillò il telefono. Molly sollevò il ricevitore mentre io
mi asciugavo le mani. «Ah-ha. No, è qui. Un attimo soltanto.» Mise il
telefono in modalità silenziosa e disse: «È Jason. Vuoi parlargli?
Sembra che stia dando di matto.»

«Certo che voglio.»

Presi la cornetta e aspettai. Molly mi guardò a lungo, poi alzò gli
occhi al cielo e uscì dalla cucina. Privacy. «Jase? Che succede?»

«Ho bisogno che tu venga qui, Tyler.» La voce era tesa, soffocata.
«Ora.»

«Hai qualche problema?»

«Sì, cazzo, ho un problema. E ho bisogno che tu venga a risolverlo.»

«È urgente?»

«Starei a chiamarti se non lo fosse?»

«Dove sei?»

«A casa.»

«Okay, ascolta, mi ci vorrà un po' di tempo se c'è molto traffico\ldots»

«Basta che vieni.»

E così dissi a Molly che avevo del lavoro urgente da sbrigare. Sorrise o
forse sogghignò e disse: «Che lavoro? Qualcuno ha perso un appuntamento?
Devi far nascere un bambino? Cosa?»

«Sono un dottore, Moll. Segreto professionale.»

«Essere un dottore non significa che tu debba essere il cagnolino di
Jason Lawton. Non devi riportargli il bastoncino ogni volta che te lo
lancia.»

«Mi dispiace interrompere la serata. Vuoi che ti dia un passaggio da
qualche parte o\ldots?»

«No, rimarrò qui finché non torni.» Mi fissava con aria di sfida, sul
piede di guerra, quasi volesse che mi opponessi.

Ma non potevo litigare. Avrebbe significato che non mi fidavo di lei. E
invece io mi fidavo di lei. Quasi sempre. «Non so di preciso quanto
tempo starò via.»

«Non importa. Mi rannicchio sul divano e guardo la tele. Se per te va
bene\ldots»

«Basta che non ti annoi.»

«Prometto che non mi annoierò.»

\separatorefiori

L'appartamento scarsamente arredato di Jason si trovava a circa trenta
chilometri di autostrada e lungo il tragitto dovetti deviare per via di
una scena del crimine, un assalto fallito a un portavalori che aveva
ucciso un'intera macchina di turisti canadesi. Jase mi fece entrare
nell'edificio e quando bussai alla porta urlò: «È aperto.»

Il grande soggiorno era più scarno di quanto non lo fosse mai stato, un
deserto di parquet in cui Jase aveva allestito il suo accampamento da
beduino. Era sdraiato sul divano. La lampada da terra, lì accanto, lo
illuminava con una luce violenta e poco lusinghiera. Era pallido e aveva
la fronte punteggiata di sudore. Gli brillavano gli occhi.

«Pensavo non potessi venire» disse. «Pensavo che forse la tua ragazza
bifolca non ti avrebbe lasciato uscire di casa.»

Gli dissi della deviazione della polizia. Poi aggiunsi: «Fammi un
favore. Non parlare di Molly in quel modo, per piacere.»

«Nel senso ``per piacere, non riferirti a lei come se parlassi di un
cafone dell'Idaho con la sensibilità di chi è cresciuto in un parcheggio
per roulotte''? Ecco qui. Per te questo e altro!»

«Ma che ti prende?»

«Domanda interessante. Le risposte possibili sono molte. Guarda.»

Si alzò in piedi.

Fu un'operazione spiacevole, debole e spasmodica. Jase era ancora alto,
ancora snello, ma la grazia fisica che una volta era sembrata così
naturale lo aveva abbandonato. Le braccia si agitavano in maniera
scomposta. Le gambe, quando finalmente riuscì a mettersi ritto,
ondeggiavano come trampoli snodati. Batteva convulsamente le palpebre.
«Ecco che mi prende» disse. Poi, la rabbia montò trasformandosi in un
altro movimento convulsivo, il suo stato emotivo instabile quanto i suoi
arti: «Guardami! C-cazzo, Tyler, \emph{guardami}.»

«Torna a sederti, Jase. Fammi controllare.» Avevo portato il kit medico.
Gli tirai su la manica e avvolsi attorno al braccio magro un misuratore
di pressione. Riuscivo a sentire la contrazione del muscolo, quasi
incontrollata.

La pressione era alta e il battito veloce. «Prendi gli anticonvulsivi?»

«Certo che prendo quegli anticonvulsivi del cazzo.»

«In orario? Nessun doppio dosaggio? Perché se ne prendi troppi, Jase, ti
fai più male che bene.»

Jason sospirò con impazienza. Poi fece qualcosa di sorprendente. Mi
afferrò un ciuffo di capelli dietro la testa, facendomi male, e tirò
verso di sé finché la mia faccia non fu vicino alla sua. Le parole
fuoriuscirono come un fiume in piena.

«Non fare il pedante con me, Tyler. Non farlo, perché ora come ora non
posso permettermelo. Forse hai dei problemi con la mia cura. Mi spiace
ma non è il momento di portare a spasso i tuoi principi del cazzo. La
posta in gioco è troppo alta. Domattina E.D. arriva in aereo per venire
alla Perielio. Pensa di avere un asso nella manica da giocare. Ci
farebbe chiudere bottega piuttosto che farmi salire sul suo trono del
cazzo. Non posso lasciare che accada, e poi guardami: ti sembro in
condizione di commettere un patricidio?» Strinse fino a quando provò
dolore - era ancora così forte - poi mollò la presa e con l'altra mano
mi allontanò spingendomi. «Perciò curami! È a questo che servi, no?»

Accostai una sedia e mi sedetti in silenzio finché non ricadde sul
divano, spossato dal suo stesso scatto d'ira. Mi vide tirare fuori dal
kit una siringa e riempirla da una bottiglietta marrone.

«Cos'è quello?»

«Sollievo temporaneo.» In realtà era un'innocua iniezione di vitamine
del gruppo B con l'aggiunta di un leggero tranquillante. Jason la guardò
con sospetto ma me la lasciò iniettare nel suo braccio. Con l'ago
fuoriuscì una minuscola goccia di sangue.

«Lo sai già cosa devo dirti» dissi. «Non esiste cura per questo
problema.»

«Nessuna cura \emph{terrestre}.»

«Che vorrebbe dire?»

«Sai cosa vuol dire.»

Parlava del processo di longevità di Wun Ngo Wen.

La ricostruzione, aveva detto Wun, curava anche una lunga lista di
disabilità genetiche. Avrebbe eliminato il ciclo della SMA dal DNA di
Jason, inibendo le proteine modificate che stavano intaccando il suo
sistema nervoso. «Ma ci vorrebbero settimane,» dissi, «e comunque, non
posso accettare l'idea che tu faccia da cavia per una procedura non
sperimentata.»

«Non è vero che non è sperimentata. I marziani lo fanno da secoli e sono
umani quanto noi. E mi dispiace, Tyler, ma non mi interessano proprio i
tuoi scrupoli professionali. Semplicemente non rientrano
nell'equazione.»

«E invece sì. Per quanto mi riguarda.»

«Allora ti chiedo: quanto ti riguarda? Se non vuoi essere coinvolto,
fatti da parte.»

«Il rischio\ldots»

«Il rischio è mio, non tuo.» Chiuse gli occhi. «Non scambiarla per
arroganza o vanità, ma è importante se vivo o muoio e anche se riesco a
camminare diritto o a pronunciare queste d-dannate consonanti. È
importante per il mondo, voglio dire. Perché mi trovo in una posizione
di particolare importanza. Non per caso. Non perché sono intelligente o
virtuoso. Sono stato prescelto. In sostanza, Tyler, sono un artefatto,
un oggetto costruito, progettato da E.D. Lawton nello stesso modo in cui
lui e tuo padre progettavano i profili alari. Sto facendo il lavoro per
cui lui mi ha costruito\ldots{} gestire la Perielio, gestire la risposta
umana allo Spin.»

«Il Presidente potrebbe non essere d'accordo. Per non parlare del
Congresso. O l'ONU, a dirla tutta.»

«Per favore. Non sto delirando. È questo il punto. Dirigere la Perielio
significa manipolare le parti interessate. Ognuna di loro. E.D. lo sa; è
decisamente cinico al riguardo. Ha trasformato la Perielio in una
gallina dalle uova d'oro per l'industria aerospaziale e ci è riuscito
creandosi amicizie e stringendo alleanze politiche nelle alte sfere.
Lusingando, implorando, esercitando pressioni e finanziando campagne di
amici. Aveva una visione, aveva gli agganci ed era nel posto giusto al
momento giusto; si è proposto col progetto degli aerostati e ha salvato
dallo Spin l'industria delle telecomunicazioni, e questo l'ha
catapultato fra le persone potenti\ldots{} e lui sa come sfruttare
un'opportunità. Senza E.D., non ci sarebbero esseri umani su Marte.
Senza E.D. Lawton, Wun Ngo Wen non esisterebbe nemmeno. Bisogna
riconoscerglielo al vecchio stronzo. È un grand'uomo.»

«Ma?»

«Ma è un uomo del suo tempo. Appartiene al periodo prima dello Spin. Le
sue motivazioni sono arcaiche. Il testimone è stato passato. O lo sarà,
se me ne occupo io.»

«Non ne capisco il senso, Jase.»

«E.D. pensa di poter ancora estorcere qualche vantaggio personale da
tutto questo. Ce l'ha con Wun Ngo Wen e odia l'idea di seminare la
galassia con i replicatori, non perché sia troppo ambiziosa ma perché
non fa bene agli affari. Il progetto di Marte ha fatto confluire
trilioni di dollari nel settore aerospaziale. Ha reso E.D. più ricco e
potente di quanto abbia mai sognato di diventare. L'ha reso famoso. E
lui continua a pensare che questa cosa conti. Crede che conti come
contava prima dello Spin, quando potevi fare politica come fosse un
gioco, una scommessa a premi. Ma la proposta di Wun non dà quel genere
di profitto. Lanciare i replicatori è un investimento banale rispetto
alla terraformazione di Marte. Possiamo farlo con un paio di razzi Delta
VII e un propulsore ionico economico. Servono soltanto una fionda e una
provetta.»

«Come può essere negativo per E.D?»

«Non serve a proteggere un settore che sta collassando. Svuota la sua
base finanziaria. E, ancora peggio, lo allontana dalla luce dei
riflettori. All'improvviso tutti rivolgeranno lo sguardo verso Wun Ngo
Wen - fra un paio di settimane ci ritroveremo in una montagna di merda
mediatica di proporzioni inaudite - e Wun mi ha scelto come leader di
questo progetto. L'ultima cosa che E.D. vuole è che il suo figlio
ingrato e un marziano grinzoso smantellino il lavoro di una vita e
lancino un'armata che costa meno di un solo aereo di linea commerciale.»

«Cosa preferirebbe fare?»

«Ha elaborato un programma su vasta scala. Il controllo dell'intero
sistema, come lo chiama lui. Cercare nuove tracce di attività degli
Ipotetici. Ispettori planetari da Mercurio a Plutone, postazioni di
ascolto sofisticate nello spazio interplanetario, missioni di sorvolo
per esplorare gli artefatti dello Spin qui e ai poli marziani.»

«È una cattiva idea?»

«Potrebbe dare qualche informazione irrilevante. Racimolare un po' di
dati e convogliare capitali liquidi nell'industria. Ecco per cosa è
progettata. Ma quello che E.D. non capisce, quello che la \emph{sua}
generazione non capisce davvero\ldots»

«Cosa, Jase?»

«È che la finestra si sta chiudendo. La finestra umana. Il nostro tempo
sulla Terra. Il tempo della Terra nell'universo. Sta per finire.
Abbiamo, credo, solo un'altra opportunità realistica di capire cosa
significa - cosa \emph{significava} - aver costruito una civiltà umana.»
Abbassò lentamente le palpebre una volta, due volte. Buona parte della
selvaggia tensione di prima lo aveva abbandonato. «Cosa significa essere
stati selezionati per questa particolare forma di estinzione. Più di
quello, comunque. Cosa significa\ldots{} cosa significa\ldots» Sollevò
lo sguardo. «Che cazzo mi hai dato, Tyler?»

«Niente di serio. Un ansiolitico leggero.»

«Una soluzione rapida?»

«Non è ciò che vuoi?»

«Credo di sì. Voglio essere presentabile per domattina, ecco ciò che
voglio.»

«Il farmaco non è una cura. È come se mi stessi chiedendo di riparare un
collegamento elettrico lasco facendoci passare più tensione. Potrebbe
funzionare per un breve periodo. Ma è incerto e carica uno stress
intollerabile su altre parti del sistema. Mi piacerebbe tanto regalarti
una bella giornata senza sintomi, ma non voglio ucciderti.»

«Se \emph{non} mi regali una giornata senza sintomi, \emph{tanto vale}
che tu mi uccida.»

«Quello che posso offrirti è il mio parere professionale.»

«E cosa posso aspettarmi dal tuo parere professionale?»

«Posso aiutarti. Credo. Un po'. Questa volta. \emph{Questa volta}, Jase.
Ma non c'è molto spazio di manovra. Devi accettarlo.»

«Nessuno di noi ha molto spazio di manovra. Tutti noi dobbiamo
accettarlo.»

Ma sospirò e sorrise quando riaprii il kit medico.

\separatorefiori

Quando tornai a casa, Molly era accovacciata sul divano e, di fronte al
televisore, guardava un recente film di successo sugli elfi, o forse
erano angeli. Tutto lo schermo emanava una luce azzurra sfocata. Quando
entrai, lo spense. Le chiesi se fosse successo qualcosa mentre ero via.

«Non molto. Hai ricevuto una telefonata.»

«Ah, sì? Chi era?»

«La sorella di Jason. Come si chiama. Diane. Quella che sta in Arizona.»

«Ha detto cosa voleva?»

«Chiacchierare e basta. E così abbiamo chiacchierato un po'.»

«Ah-ha. Di cosa avete parlato?»

Molly si voltò solo in parte, mostrandomi il profilo sullo sfondo della
luce fioca proveniente dalla camera da letto. «Di te.»

«Di qualcosa in particolare?»

«Sì. Le ho detto di smettere di chiamarti perché ora sei fidanzato con
un'altra. Le ho detto che d'ora in poi mi sarei occupata io delle tue
telefonate.»

Sgranai gli occhi.

Molly mostrò i denti in quello che io interpretai come un sorriso. «E
dai, Tyler, impara a stare agli scherzi. Le ho detto che eri fuori. Va
bene così?»

«Le hai detto che ero fuori?»

«Sì, le ho detto che eri fuori. Non ho detto dove. Perché effettivamente
non me l'hai detto.»

«Ha specificato se era urgente?»

«Non mi è sembrato urgente. Richiamala se vuoi. Fai pure\ldots{} non c'è
problema.»

Ma mi stava mettendo di nuovo alla prova. «Non c'è fretta» risposi.

«Bene.» Le vennero le fossette sulle guance. «Perché ho altri
programmi.»
